% V033 图像先行：不重复展开证明中的四种作法
\begin{explanationbox}
\textbf{\textcolor{CExplanation}{先判断交点位置，再决定面积加减}}
\begin{center}
\begin{tikzpicture}[x=1cm,y=1cm,font=\small,>=Stealth,
  every node/.style={inner sep=3pt},line cap=round,line join=round]
\path[use as bounding box] (-1,-1.6) rectangle (11,5.5);
\node[font=\small\bfseries,text=CStatement] at (5,5.1) {辅助点在反向延长线：大面积减去小面积};

\coordinate (E) at (1.166667,0); \coordinate (O) at (2.5,1);
\coordinate (A) at (6.5,4); \coordinate (B) at (4.5,0);
\fill[CStatement!17] (E)--(A)--(B)--cycle;
\fill[CTrap!25] (E)--(O)--(B)--cycle;
\draw[thick] (E)--(A)--(B)--cycle;
\draw[very thick,CStatement] (O)--(B);
\draw[very thick,CTrap] (E)--(B);
\draw[dashed] (O)--(7.4,1) (A)--(7.4,4) (B)--(7.4,0);
\draw[<->] (7.1,0)--node[right] {$4$}(7.1,4);
\draw[<->] (5.2,0)--node[right] {$1$}(5.2,1);
\node[below] at (E) {$E_h(-\frac43,-1)$};
\node[left] at (O) {$O(0,0)$};
\node[above] at (A) {$A(4,3)$}; \node[below right] at (B) {$B(2,-1)$};
\node at (8.7,4.1) {蓝片：$\triangle OAB$};
\node[align=center] at (8.7,3.1) {减去红片\\$\triangle E_hOB$};

\node[align=center] at (5,-1.15) {$BE_h=\frac{10}{3}$，$S=\frac12\cdot\frac{10}{3}\cdot(4-1)=5$。};
\end{tikzpicture}
\end{center}

图中 $E_h,O,A$ 依次共线。大三角形 $E_hAB$ 由原三角形 $OAB$ 与红色小三角形 $E_hOB$ 拼成，故必须相减。
同一例中，竖直辅助点为 $E_v=(2,3/2)$，位于线段 $OA$ 内，改用竖直法则把两片相加。
\medskip

\textbf{\textcolor{CExplanation}{四种初中作法：起点与方向都可选择}}
\begin{itemize}[leftmargin=*,itemsep=0.4em]
\item 过 $B$ 作竖直线：求 $y_{E_v}=x_2y_1/x_1$，合并水平高为 $|x_1|$。
\item 过 $B$ 作水平线：求 $x_{E_h}=x_1y_2/y_1$，合并竖直高为 $|y_1|$。
\item 过 $A$ 作竖直线交 $OB$ 于 $F_v$：求 $y_{F_v}=x_1y_2/x_2$，合并水平高为 $|x_2|$。
\item 过 $A$ 作水平线交 $OB$ 于 $F_h$：求 $x_{F_h}=x_2y_1/y_2$，合并竖直高为 $|y_2|$。
\end{itemize}
相似比的分母被底乘高约去，得到“交叉乘减”。分母为零时，直接以相应的水平或竖直边 $OA$ 或 $OB$ 为底；四幅构造图与独立推导见方法二至五。
\end{explanationbox}
\clearpage
\begin{explanationbox}
\textbf{\textcolor{CExplanation}{应用迁移：C088 双曲线渐近线}}
\begin{center}
\begin{tikzpicture}[x=1cm,y=1cm,font=\small,>=Stealth,
  every node/.style={inner sep=3pt},line cap=round,line join=round]
\path[use as bounding box] (-1,-1.6) rectangle (11,5.5);
\node[font=\small\bfseries,text=CStatement] at (5,5.1) {两条渐近线上的点：斜率一正一负};

\begin{scope}[shift={(4,1.7)},scale=0.65]
\draw[->,gray] (-2,0)--(7.5,0) node[right] {$x$};
\draw[->,gray] (0,-3)--(0,4.2) node[above] {$y$};
\draw[CStatement,thick] (-2,-1.3333)--(6.5,4.3333);
\draw[CExample,thick,dashed] (-2,1.3333)--(4,-2.6667);
\fill[CStatement!17] (0,0)--(6,4)--(3,-2)--cycle;
\draw[thick] (6,4)--(3,-2);
\node[above left] at (0,0) {$O$};
\node[right] at (6,4) {$A(6,4)$}; \node[right] at (3,-2) {$B(3,-2)$};
\node[left] at (-3,2) {$y=-\frac23x$};
\node[above left] at (2.8,2.6) {$y=\frac23x$};
\end{scope}

\node[align=center] at (5,-1.15) {$S=\frac12|6(-2)-3\cdot4|=12=\frac23|6\cdot3|$。};
\end{tikzpicture}
\end{center}

一般地，设 $a,b>0$，两渐近线为 $y=\pm(b/a)x$，
取 $A=(x_1,\frac ba x_1)$、$B=(x_2,-\frac ba x_2)$，则
\[
S=\frac12\left|x_1\left(-\frac ba x_2\right)-x_2\left(\frac ba x_1\right)\right|
=\boxed{\frac ba|x_1x_2|}.
\]
两个交叉乘积符号相反，做差产生因子 $2$，抵消面积公式中的 $1/2$。
\end{explanationbox}
